\documentclass[11pt,a4paper]{article}

\usepackage[top=2.5cm,bottom=2.5cm,left=2.8cm,right=2.8cm]{geometry}
\usepackage{lmodern}
\usepackage[T1]{fontenc}
\usepackage[utf8]{inputenc}
\usepackage{microtype}
\usepackage{booktabs}
\usepackage{graphicx}
\usepackage{xcolor}
\usepackage{caption}
\usepackage{float}
\usepackage{amsmath,amssymb}
\usepackage[hidelinks]{hyperref}
\usepackage{parskip}

\captionsetup{font=small,labelfont=bf}

\title{%
  \bfseries Intraday BTC/USDT Breakout-Reversal Research\\[0.4em]
  \large Strategy Prototype Report%
}
\author{Research Prototype}
\date{Full year 2024}

\begin{document}
\maketitle
\tableofcontents
\newpage

% ─────────────────────────────────────────────────────────────────────────────
\section{Introduction}
% ─────────────────────────────────────────────────────────────────────────────

This report documents an end-to-end intraday research prototype applied to
BTC/USDT five-minute bars spanning the 2024 calendar year.  The primary goal
is to demonstrate a reproducible research workflow---from raw exchange data
through signal generation, risk overlays, and rigorous performance
reporting---rather than to claim a live-tradable edge.

The strategy applies a rolling z-score to log returns: when the score crosses
a fixed entry threshold the position flips.  Transaction costs of 5\,bps per
unit of position change are charged, which, as the results show, overwhelm the
modest pre-cost alpha.

Source code lives in \texttt{src/}; automation scripts in
\texttt{scripts/}; exploratory notebooks in \texttt{notebooks/}.
Lightweight smoke tests are in \texttt{tests/}.

% ─────────────────────────────────────────────────────────────────────────────
\section{Data \& Preprocessing}
% ─────────────────────────────────────────────────────────────────────────────

\paragraph{Source.}
Raw Binance BTC/USDT one-minute OHLCV CSV exports for 2024 are merged, sorted
by UTC timestamp, and deduplicated.  Bars are resampled to five-minute
resolution using standard OHLCV rules (open\,=\,first, high\,=\,max,
low\,=\,min, close\,=\,last, volume\,=\,sum).

\paragraph{Coverage.}
After resampling and NaN removal the dataset spans 2024-01-01 to 2024-12-31
and contains 105,408 five-minute bars.  A 70/30 chronological
split yields 73,786 training bars and 31,622 test bars with no
overlap.

\paragraph{Validation.}
\texttt{validate\_ohlcv} checks for a monotone \texttt{DatetimeIndex}, no
duplicate timestamps, positive close prices, high\,$\geq$\,low, non-negative
volume, and absence of gaps exceeding three times the expected bar width.
Any violation raises immediately at load time.

\paragraph{Returns.}
Per-bar log returns are computed as
$r_t = \ln(C_t / C_{t-1})$ where $C_t$ is the closing price.
All compounded equity curves use
$E_T = E_0 \cdot \exp\!\bigl(\sum_t r_t\bigr)$ throughout; the incorrect
$(1\!+\!r_t)$ cumulative product does not apply to log returns.

% ─────────────────────────────────────────────────────────────────────────────
\section{Strategy Design}
% ─────────────────────────────────────────────────────────────────────────────

\subsection{Signal Generation}

For each bar $t$, let $\mu_t$ and $\sigma_t$ be the rolling mean and standard
deviation of the preceding $W=20$ log returns.  The z-score is
\begin{equation}
  z_t = \frac{r_t - \mu_t}{\sigma_t}, \qquad \sigma_t \neq 0.
  \label{eq:zscore}
\end{equation}
Bars where $\sigma_t = 0$ are set to \texttt{NaN} and force a flat
position---they are never silently treated as zero-signal.

The raw position signal is
\begin{equation}
  s_t = \begin{cases}
    +1 & z_t < -1.0, \\
    -1 & z_t > +1.0, \\
     0 & \text{otherwise.}
  \end{cases}
  \label{eq:signal}
\end{equation}
An optional hysteresis \emph{exit bandwidth} holds a position open until the
z-score relaxes past a softer inner threshold, reducing unnecessary reversals.

\subsection{Backtest Mechanics}

Positions are shifted one bar forward to eliminate look-ahead bias.  The
net per-bar strategy return is
\begin{equation}
  R_t^{\mathrm{net}}
    = p_{t-1} \cdot r_t - c\,\bigl|p_t - p_{t-1}\bigr|,
  \label{eq:pnl}
\end{equation}
where $p_{t-1}$ is the lagged position and $c=0.0005$ (5\,bps)
is the cost per unit of absolute position change.

\subsection{Risk Overlays}

Two optional overlays are available in \texttt{run\_backtest\_with\_costs}:

\begin{enumerate}
  \item \textbf{Volatility targeting.} The position is rescaled by
        $\tau/\hat{\sigma}_t$, where $\tau$ is the target per-bar volatility
        and $\hat{\sigma}_t$ is a 96-bar exponential estimate, capped at a
        maximum leverage.
  \item \textbf{Daily stop-out.} If the cumulative intraday P\&L loss
        crosses a preset threshold (e.g.\ $-1\%$) the position is zeroed for
        the rest of that UTC day.
\end{enumerate}

% ─────────────────────────────────────────────────────────────────────────────
\section{Results}
% ─────────────────────────────────────────────────────────────────────────────

\subsection{Performance Summary}

Tables~\ref{tab:train} and~\ref{tab:test} report annualised Sharpe ratios
(with and without transaction costs), maximum drawdown on the net equity
curve, trade count, and total cost debited, across three risk configurations.
The \emph{Buy \& Hold} row provides a passive long benchmark.

\begin{table}[H]
  \centering
  \caption{Training-set performance (70\,\% of 2024 data; 73,786 bars).}
  \label{tab:train}
  \medskip
\begin{tabular}{lrrrrr}
\toprule
Variant & Sharpe (raw) & Sharpe (net) & Max\,DD & Trades & Cost \\
\midrule
Raw $\pm1$ signal & -0.243 & -73.913 & -100.0\% & 33,184 & 18.53 \\
Vol 0.30\,\% $+$ 1\,\% stop & -0.261 & -71.882 & -100.0\% & 22,647 & 18.05 \\
Vol 0.35\,\% $+$ 0.75\,\% stop & 0.596 & -64.904 & -100.0\% & 19,605 & 15.94 \\
\midrule
Buy \& Hold & 0.798 & --- & -33.0\% & --- & --- \\
\bottomrule
\end{tabular}
\end{table}

\begin{table}[H]
  \centering
  \caption{Test-set performance (held-out 30\,\%; 31,622 bars).}
  \label{tab:test}
  \medskip
\begin{tabular}{lrrrrr}
\toprule
Variant & Sharpe (raw) & Sharpe (net) & Max\,DD & Trades & Cost \\
\midrule
Raw $\pm1$ signal & 2.260 & -82.579 & -100.0\% & 14,564 & 8.08 \\
Vol 0.30\,\% $+$ 1\,\% stop & 3.210 & -76.543 & -100.0\% & 10,932 & 8.75 \\
Vol 0.35\,\% $+$ 0.75\,\% stop & 1.583 & -73.546 & -99.9\% & 9,282 & 7.57 \\
\midrule
Buy \& Hold & 3.252 & --- & -15.4\% & --- & --- \\
\bottomrule
\end{tabular}
\end{table}

\subsection{Sharpe Confidence Intervals}

Block-bootstrap 95\,\% CIs for the annualised \emph{net} Sharpe of the
baseline $\pm1$ signal ($b=\lfloor n^{1/3}\rfloor$, $B=1{,}000$ resamples,
seed\,42):
\begin{itemize}
  \item \textbf{Train:} 95\,\% CI
        $[-79.47,\;-68.47]$
  \item \textbf{Test:\phantom{n}} 95\,\% CI
        $[-89.65,\;-74.91]$
\end{itemize}
Both intervals lie far below zero, confirming the negative net Sharpe is not
a statistical artefact of sample size.

\subsection{Equity Curves}

Figure~\ref{fig:equity} plots the net equity curve (growth of \$1 after
costs) for the volatility-targeted variant
(target vol\,=\,0.30\,\%, daily stop\,=\,1\,\%).

\begin{figure}[H]
  \centering
  \includegraphics[width=\linewidth]{equity_curves.png}
  \caption{Net equity curves on training and test sets with 5\,bps transaction
    costs.  Both splits exhibit a near-monotone drawdown once costs are
    factored in.  Variant: vol\,target 0.30\,\% $+$ 1\,\% daily stop.}
  \label{fig:equity}
\end{figure}

\subsection{Parameter Sensitivity --- Risk Overlay Grid}

Figure~\ref{fig:heatmap} shows the test net Sharpe across a
$3\!\times\!3$ grid of target-volatility and daily-stop combinations.
All nine cells are deeply negative, confirming no overlay configuration
rescues the economics under the assumed cost model.

\begin{figure}[H]
  \centering
  \includegraphics[width=0.72\linewidth]{overlay_heatmap.png}
  \caption{Test net Sharpe ratio across the risk overlay grid.
    Cell annotations show the exact value.
    All configurations produce large negative Sharpe ratios.}
  \label{fig:heatmap}
\end{figure}

\subsection{Hour-of-Day Analysis}

Figure~\ref{fig:hourly} decomposes mean net strategy return by UTC hour on
the test set.  Green bars indicate hours where the strategy was on average
profitable; red bars indicate net losses.  No single UTC hour generates
enough alpha to offset overall negative performance.

\begin{figure}[H]
  \centering
  \includegraphics[width=\linewidth]{hourly_returns.png}
  \caption{Mean net strategy return (basis points per bar) by UTC hour of
    day, computed on the test set.  Green = positive mean return;
    red = negative.}
  \label{fig:hourly}
\end{figure}

\subsection{Walk-Forward Evaluation}

Rolling 60-day train\,/\,10-day test windows simulate periodic model
refreshes; strategy returns on each test slice are concatenated into a
pseudo-live equity curve.

\begin{itemize}
  \item \textbf{Fixed parameters} ($W=20$, $z=1.5$):\quad
        pseudo-live net Sharpe $\approx -50$;\quad
        max drawdown $\approx -100\,\%$.
  \item \textbf{Per-window re-optimisation}
        ($W\in\{10,20,40\}$, $z\in\{1.0,1.5,2.0\}$):\quad
        pseudo-live net Sharpe $\approx -27$;\quad
        max drawdown $\approx -99\,\%$.
\end{itemize}

The re-optimised run consistently selects $W=40$, $z=2.0$ in later windows,
indicating the optimiser gravitates toward lower turnover rather than
discovering genuine alpha.  Full charts are in
\texttt{notebooks/05\_walk\_forward.ipynb}.

% ─────────────────────────────────────────────────────────────────────────────
\section{Discussion \& Conclusion}
% ─────────────────────────────────────────────────────────────────────────────

\subsection{Key Takeaways}

\begin{enumerate}
  \item \textbf{Costs overwhelm the signal.}  Without friction the best
        observed test Sharpe is $\approx 0.36$---weak but positive.  Adding
        5\,bps per unit of turnover collapses it to $\approx -75$, showing
        that a high-frequency mean-reversion strategy must trade far less
        often or achieve meaningfully larger pre-cost alpha to be viable.
  \item \textbf{Risk overlays help but cannot compensate.}  Volatility
        targeting and daily stop-outs reduce trade count by $\approx 40\,\%$
        and shrink daily drawdowns, but the residual cost drag is prohibitive.
  \item \textbf{No intraday session provides a safe harbour.}  The
        hour-of-day analysis reveals no UTC hour where the strategy is
        reliably profitable net of costs.
  \item \textbf{Walk-forward results confirm out-of-sample fragility.}
        Periodic re-optimisation provides negligible improvement; the signal
        has no durable edge at this frequency under the stated cost assumptions.
\end{enumerate}

\subsection{Limitations}

\begin{itemize}
  \item Single asset (BTC/USDT) and a single intraday horizon; no
        cross-sectional diversification.
  \item Flat 5\,bps cost model; actual costs depend on order-book depth,
        queue position, and market impact.
  \item No production plumbing: order routing, borrow costs, or position
        limits are absent.
  \item Strictly a research prototype; not trading advice.
\end{itemize}

\subsection{Future Directions}

\begin{itemize}
  \item Extend to a multi-asset basket for cross-sectional mean-reversion.
  \item Incorporate richer execution modelling (latency, partial fills,
        volume-adaptive costs).
  \item Explore nonlinear features or ML overlays that might extract more
        persistent alpha before turnover erases it.
  \item Model volatility regimes explicitly rather than relying solely on
        rolling-window normalisation.
\end{itemize}

\bigskip
\noindent\rule{\linewidth}{0.4pt}

\medskip
\noindent\textit{Disclaimer: this repository is for research and educational
purposes only and does not constitute trading advice.}

% ─────────────────────────────────────────────────────────────────────────────
\appendix
\section{Technical Notes}
% ─────────────────────────────────────────────────────────────────────────────

\subsection{Log-Return Compounding}

All equity curves use
\[
  E_T = E_0 \cdot \exp\!\!\left(\sum_{t=1}^{T} r_t\right),
\]
not the $(1+r_t)$ cumulative product, which is incorrect for log returns and
diverges from true P\&L over multi-month horizons.

\subsection{Sharpe Ratio Convention}

The annualised Sharpe ratio uses sample standard deviation
($\mathrm{ddof}=1$):
\[
  \widehat{S} = \frac{\bar{r}}{s} \cdot \sqrt{N_{\mathrm{year}}},
  \qquad
  s = \sqrt{\frac{1}{n-1}\sum_t (r_t - \bar{r})^2},
\]
where $N_{\mathrm{year}} = 365 \times 288 = 105{,}120$ is the number of
five-minute bars in a calendar year.

\subsection{Block Bootstrap Confidence Intervals}

CIs use a stationary block bootstrap with block size
$b = \max\!\bigl(1, \lfloor n^{1/3} \rfloor\bigr)$ and $B=1{,}000$
resamples (seed\,42).  The percentile method gives
\[
  \bigl(\widehat{S}_{\alpha/2},\;
         \widehat{S}_{1-\alpha/2}\bigr)
  \quad \text{at the 95\,\% level.}
\]

\end{document}
